\begin{table}
\caption{CV $\times$ sex $\times$ SES triple-interaction coefficient (age-18 cutoff).}
\label{tab:triple_T18}
\begin{tabular}{lrrrrlrr}
\toprule
SES measure & Coef. (pp) & SE (pp) & z & p & BW (days) & p10 & p90 \\
\midrule
NBI & -26.251 & 19.474 & -1.348 & 0.178 & 47.4 & 0.026 & 0.134 \\
PCA index & -1.100 & 0.679 & -1.620 & 0.105 & 47.4 & 1.599 & 5.109 \\
\bottomrule
\end{tabular}
\par\vspace{0.6em}
{\footnotesize\begin{minipage}{\linewidth}\emph{Note:} Coefficient on $T \times$ female $\times$ SES from Equation~\ref{eq:triple}, in percentage points per one-unit change in the SES measure. The two measures are on different scales --- NBI is a proportion $(0,1)$, the PCA index is in index units --- so the columns p10 and p90 give the range each actually spans: the implied change in the gender gap across the observed distribution is Coef.\,$\times$\,(p90 $-$ p10). Standard errors clustered by electoral circuit (CRV1).\end{minipage}}
\end{table}
